\SourcePage{174}
\SGAsection{II}{等变对应}
\SGAsection{5}{非交换迹}
本节固定一个\Term{topos}{拓扑斯} $T$ 及其中一个交换含幺\Term{ring}{环} $K$。所考虑的\Term{ring}{环}均假定为结合含幺\Term{ring}{环}，但未必交换。所谓 $K$-\Term{algebra}{代数}，是指 $T$ 中一个（结合、含幺）\Term{ring}{环} $A$，配备像包含于\Term{central}{中心的}（含幺）\Term{homomorphism}{同态} $K\to A$。

\SGAsection{5.0}{引言}
设 $E$ 是一个 $K$-\Term{module}{模}，\Term{local-scheme}{局部}为某个\Term{finite-type}{有限型}\Term{free-module}{自由模}的\Term{direct-summand}{直和因子}。这时有一个（具有\Term{local-nature}{局部性质}的）\Term{trace}{迹}\Term{homomorphism}{同态}
\begin{equation}\operatorname{Tr}_K:\mathcal{H}\!om_K(E,E)\longrightarrow K,\tag{5.0.1}\end{equation}
它由“\Term{evaluation-map}{求值箭头}”
\begin{equation}\mathcal{H}\!om_K(E,K)\otimes_K E\longrightarrow K,\qquad f\otimes x\longmapsto f(x),\tag{5.0.2}\end{equation}
与\Term{canonical}{典范}“\Term{tensor-product}{张量积}”\Term{isomorphism}{同构}
\begin{equation}\mathcal{H}\!om_K(E,K)\otimes_K E\longrightarrow\mathcal{H}\!om_K(E,E),\qquad f\otimes y\longmapsto(x\longmapsto f(x)y)\tag{5.0.3}\end{equation}
的逆复合而成。更一般地，\Term{arrow}{箭头} 5.0.2 和 5.0.3 对任意 $K$-\Term{module}{模} $E$ 均有定义；于是对任意 $K$-\Term{module}{模}\Term{complex}{复形} $E$，可得\Term{complex}{复形}的\Term{homomorphism}{同态}
\begin{equation}\mathcal{H}\!om_K^\bullet(E,K)\otimes_K E\longrightarrow K,\qquad f\otimes x\longmapsto f(x),\tag{5.0.4}\end{equation}
\begin{equation}\mathcal{H}\!om_K^\bullet(E,K)\otimes_K E\longrightarrow\mathcal{H}\!om_K^\bullet(E,E),\qquad f\otimes y\longmapsto(x\longmapsto f(x)y).\tag{5.0.5}\end{equation}
对 5.0.4、5.0.5 进行\Term{derive}{导出}，对 $E\in\operatorname{ob}D^-(K)$ 得到 $D(K)$ 中的\Term{arrow}{箭头}：
\begin{equation}R\mathcal{H}\!om_K(E,K)\overset L\otimes_K E\longrightarrow K,\tag{5.0.6}\end{equation}
\begin{equation}R\mathcal{H}\!om_K(E;K)\overset L\otimes_K E\longrightarrow R\mathcal{H}\!om_K(E,E).\tag{5.0.7}\end{equation}

\SourcePage{175}
若 $E$ 是\Term{perfect-complex}{完美的}（SGA 6 I），则 5.0.7 是\Term{isomorphism}{同构}（归结到 $E=K$ 的情形即可容易地看出）。将 5.0.7 的逆与 5.0.6 复合，便得到一个具有整体性质的\Term{trace}{迹}\Term{homomorphism}{同态}
\begin{equation}\operatorname{Tr}_K:R\mathcal{H}\!om_K(E,E)\longrightarrow K,\tag{5.0.8}\end{equation}
它推广了 5.0.1（更详细的说明见 SGA 6 I 7、8）。

这里我们想考察，在上述构造中能在多大程度上以一个 $K$-\Term{algebra}{代数} $A$ 代替 $K$（我们主要关心 $A$ 是一个普通\Term{finite-group}{有限群}的 $K$-\Term{group-algebra}{群代数}的情形），特别是定义一个类似于 5.0.8 的\Term{trace}{迹}\Term{arrow}{箭头}，推广 Stallings 在 [5] 中引入、Grothendieck 在本讨论班第 XI 讲中发展的“\Term{noncommutative-trace}{非交换迹}”概念（如引言所述，该讲由 I. Bucur 撰写的讲稿不幸遗失；不过读者可参阅 Deligne 的讲稿 SGA $4\frac12$ Rapport，其中给出了主要论点）。在 5.0.3 中将 $K$ 换成 $A$ 并无困难，但在 5.0.2 中情况就不同了：若 $E$ 是左 $A$-\Term{module}{模}，则\Term{evaluation-map}{求值箭头}
\[\mathcal{H}\!om_A(E,A)\otimes_K E\longrightarrow A,\qquad f\otimes x\longmapsto f(x)\]
一般不能通过 $\mathcal{H}\!om_A(E,A)\otimes_A E$ \Term{resolution}{分解}，但取商后给出一个\Term{arrow}{箭头}
\begin{equation}\mathcal{H}\!om_A(E,A)\otimes_A E\longrightarrow A_\natural,\tag{5.0.9}\end{equation}
其中 $A_\natural$ 是 $A$ 对由形如 $ab-ba$ 的\Term{local-section}{局部截面}\Term{local-scheme}{局部}生成的 $K$-\Term{submodule}{子模}的商 $K$-\Term{module}{模}。若 $E$ \Term{local-scheme}{局部}为某个\Term{finite-type}{有限型}自由左 $A$-\Term{module}{模}的\Term{direct-summand}{直和因子}，则类似于 5.0.3 的\Term{tensor-product}{张量积}\Term{arrow}{箭头}
\begin{equation}\mathcal{H}\!om_A(E,A)\otimes_A E\longrightarrow\mathcal{H}\!om_A(E,E)\tag{5.0.10}\end{equation}
是\Term{isomorphism}{同构}；将 5.0.10 的逆与 5.0.9 复合，就得到\Term{trace}{迹}\Term{homomorphism}{同态}
\SourcePage{176}
\begin{equation}\operatorname{Tr}_A:\mathcal{H}\!om_A(E,E)\longrightarrow A_\natural,\tag{5.0.11}\end{equation}
当 $A=K$ 时，它化为 5.0.1。另外容易验证（见 5.7），当 $T$ 是点\Term{topos}{拓扑斯}时，5.0.11 与 Stallings 和 Grothendieck 的\Term{trace}{迹}\Term{homomorphism}{同态}一致。本节将考察\Term{arrow}{箭头} 5.0.9、5.0.10、5.0.11 向\Term{complex}{复形}的扩展及其\Term{derive}{导出}，并研究所得\Term{arrow}{箭头}对于\Term{extension-of-scalars}{标量扩张}、\Term{restriction-of-scalars}{标量限制}及某些\Term{external-tensor-product}{外张量积}的行为。实际上，我们将\Term{derive}{导出}比 5.0.9 和 5.0.10 略为一般的\Term{arrow}{箭头}，使这些构造能用于第 6 节“非交换” Lefschetz--Verdier \Term{local-term}{局部项}的定义。

\SGAsection{5.1}{记号}
若 $A$ 是一个 $K$-\Term{algebra}{代数}，以 $A^\circ$ 表示 $A$ 的\Term{opposite-algebra}{反代数}，以 $A^e=A\otimes_K A^\circ$ 表示 $A$ 的\Term{enveloping-algebra}{包络代数}。给定一个 $A$-\Term{bimodule}{双模}、一个左 $A^e$-\Term{module}{模}，或一个右 $A^e$-\Term{module}{模}，是同一回事。特别地，$A$ 是一个 $A$-\Term{bimodule}{双模}，因而既是左 $A^e$-\Term{module}{模}也是右 $A^e$-\Term{module}{模}。若 $E$ 是一个\Term{bimodule}{双模}，回顾（[1] IX 4），
\[H_0(A,E)=A\otimes_{A^e}E=E\otimes_{A^e}A\]
是 $E$ 对由形如 $ax-xa$ 的\Term{local-section}{局部截面}\Term{local-scheme}{局部}生成的 $K$-\Term{submodule}{子模}的商 $K$-\Term{module}{模}，其中 $x\in E$、$a\in A$\footnote{若 $M$ 是一个\Term{sheaf}{层}，记号 $x\in M$ 表示 $x$ 是 $M$ 的一个\Term{local-section}{局部截面}。}。有时以 $E_\natural$ 代替 $H_0(A,E)$。特别地，$A_\natural$ 就是 5.0.9 中定义的 $K$-\Term{module}{模}。若 $A=K[G]$，其中 $G$ 为普通\Term{finite-group}{有限群}，则 $A_\natural$ 正是以 $G$ 的\Term{conjugacy-class}{共轭类}\Term{set}{集合} $G_\natural$ 为基的自由 $K$-\Term{module}{模}。

若 $A,B$ 是 $K$-\Term{algebra}{代数}，我们用 Cartan--Eilenberg 的记号 ${}_AE$、$E_B$、${}_AE_B$ 分别表示左 $A$-\Term{module}{模}、右 $B$-\Term{module}{模}及\Term{left-action}{左作用}为 $A$、\Term{right-action}{右作用}为 $B$ 的\Term{bimodule}{双模}。最后，用 $D(A)$（相应地 $D(A,B)$）表示左 $A$-\Term{module}{模}\Term{category}{范畴}（相应地\Term{left-action}{左作用}为 $A$、\Term{right-action}{右作用}为 $B$ 的 $A$-$B$-\Term{bimodule}{双模}\Term{category}{范畴}）的\Term{derived-category}{导出范畴}。

\SourcePage{177}
\textbf{5.2.} 设 $A,B$ 是 $K$-\Term{algebra}{代数}。回顾，在 $({}_AE,{}_BF_A,{}_BG)$ 的情形有\Term{canonical}{典范}\Term{isomorphism}{同构}
\begin{equation}\mathcal{H}\!om_A(E,\mathcal{H}\!om_B(F,G))\xrightarrow\sim\mathcal{H}\!om_B(F\otimes_A E,G),\qquad f\longmapsto\bigl(y\otimes x\longmapsto f(x)(y)\bigr).\tag{5.2.1}\end{equation}
它\Term{derive}{导出}为\Term{isomorphism}{同构}
\begin{equation}R\mathcal{H}\!om_A(E,R\mathcal{H}\!om_B(F,G))\xrightarrow\sim R\mathcal{H}\!om_B(F\overset L\otimes_A E,G),\tag{5.2.2}\end{equation}
其中 $E\in\operatorname{ob}D^-(A)$、$F\in\operatorname{ob}D(B,A)$、$G\in\operatorname{ob}D^+(B)$。事实上，可假定 $E$ \Term{bounded-above}{上有界}且各\Term{set-component}{分量}\Term{flat}{平坦}，$G$ \Term{bounded-below}{下有界}且各\Term{set-component}{分量}\Term{injective}{内射}。这时有
\[R\mathcal{H}\!om_B(F\overset L\otimes_A E,G)=\mathcal{H}\!om_B^\bullet(F\otimes_A E,G),\qquad R\mathcal{H}\!om_B(F,G)=\mathcal{H}\!om_B^\bullet(F,G),\]
以及
\[R\mathcal{H}\!om_A(E,\mathcal{H}\!om_B^\bullet(F,G))=\mathcal{H}\!om_A^\bullet(E,\mathcal{H}\!om_B^\bullet(F,G)),\]
因为若 $E$ \Term{acyclic}{无环}，则 $F\otimes_A E$ \Term{acyclic}{无环}，所以 $\mathcal{H}\!om_A^\bullet(E,\mathcal{H}\!om_B^\bullet(F,G))\simeq\mathcal{H}\!om_B^\bullet(F\otimes_A E,G)$ \Term{acyclic}{无环}。对如上选取的 $E,G$，5.2.2 由 5.2.1 定义的\Term{complex}{复形}\Term{isomorphism}{同构}
\[\mathcal{H}\!om_A^\bullet(E,\mathcal{H}\!om_B^\bullet(F,G))\xrightarrow\sim\mathcal{H}\!om_B^\bullet(F\otimes_A E,G)\]
给出。注意，若 $A,B$ 在 $K$ 上\Term{flat}{平坦}，则 5.2.2 在条件 $E\in\operatorname{ob}D(A)$、$F\in\operatorname{ob}D^-(B,A)$、$G\in\operatorname{ob}D^+(B)$ 下仍然成立；以\Term{flat}{平坦} $(B\otimes_K A^\circ)$-\Term{module}{模}\Term{resolution}{消解} $F$ 即可看出。由 5.2.2 得到\Term{isomorphism}{同构}
\begin{equation}\operatorname{Hom}_A(E,R\mathcal{H}\!om_B(F,G))\xrightarrow\sim\operatorname{Hom}_B(F\overset L\otimes_A E,G).\tag{5.2.3}\end{equation}

\textbf{5.3.} 设 $A$ 是一个 $K$-\Term{algebra}{代数}。在 $(E_A,{}_AF)$ 的情形，有\Term{canonical}{典范}\Term{isomorphism}{同构}
\begin{equation}E\otimes_A F\xrightarrow\sim(E\otimes_K F)\otimes_{A^e}A,\qquad x\otimes y\longmapsto(x\otimes y)\otimes1.\tag{5.3.1}\end{equation}
另一方面，在 $({}_AE,{}_AF_A)$ 的情形，\Term{evaluation-map}{求值箭头}
\begin{equation}\mathcal{H}\!om_A(E,F)\otimes_K E\longrightarrow F,\qquad f\otimes x\longmapsto f(x)\tag{5.3.2}\end{equation}
是 $A^e$-\Term{linear}{线性}的，因而通过 5.3.1，应用 $-\otimes_{A^e}A$ 后诱导一个\Term{arrow}{箭头}
\begin{equation}\mathcal{H}\!om_A(E,F)\otimes_A E\longrightarrow F\otimes_{A^e}A\ (=F_\natural).\tag{5.3.3}\end{equation}
\SourcePage{178}
换言之，5.3.2 与\Term{canonical}{典范}\Term{arrow}{箭头} $F\to F_\natural$、$x\mapsto x\otimes1$ 的复合（唯一地）通过 5.3.3 \Term{resolution}{分解}。

假设 $A$ 在 $K$ 上\Term{flat}{平坦}。则 5.3.1 \Term{derive}{导出}为\Term{isomorphism}{同构}
\begin{equation}E\overset L\otimes_A F\xrightarrow\sim(E\overset L\otimes_K F)\overset L\otimes_{A^e}A,\tag{5.3.4}\end{equation}
其中 $E\in\operatorname{ob}D^-(A^\circ)$、$F\in\operatorname{ob}D^-(A)$（用\Term{flat}{平坦}\Term{module}{模}\Term{resolution}{消解} $E,F$）。另一方面，5.3.3 \Term{derive}{导出}为 $D(K)$ 中一个称为\Term{evaluation-map}{求值箭头}的\Term{arrow}{箭头}
\begin{equation}R\mathcal{H}\!om_A(E,F)\overset L\otimes_A E\longrightarrow F\overset L\otimes_{A^e}A,\tag{5.3.5}\end{equation}
其中 $E\in\operatorname{ob}D^-(A)$、$F\in\operatorname{ob}D^b(A^e)$，且 $R\mathcal{H}\!om_A(E,F)\in\operatorname{ob}D^b(A^\circ)$。先以\Term{flat}{平坦}\Term{module}{模}\Term{resolution}{消解} $E$、以\Term{injective}{内射} $A^e$-\Term{module}{模}\Term{resolution}{消解} $F$，将 5.3.2 \Term{derive}{导出}为 $D^-((A^e)^\circ)$ 中的\Term{arrow}{箭头}
\begin{equation}R\mathcal{H}\!om_A(E,F)\overset L\otimes_K E\longrightarrow F,\tag{5.3.6}\end{equation}
然后对 5.3.6 应用\Term{functor}{函子} $-\overset L\otimes_{A^e}A$，并利用 5.3.4。

\textbf{5.4.} 设 $A$ 是一个 $K$-\Term{algebra}{代数}。在 $({}_AE,{}_AF_A,{}_AG)$ 的情形有一个\Term{canonical}{典范}\Term{arrow}{箭头}
\begin{equation}\mathcal{H}\!om_A(E,F)\otimes_A G\longrightarrow\mathcal{H}\!om_A(E,F\otimes_A G),\qquad u\otimes y\longmapsto(x\longmapsto u(x)\otimes y),\tag{5.4.1}\end{equation}
它通过\Term{adjunction}{伴随}（5.2.1）对应于由\Term{evaluation-map}{求值箭头} $E\otimes_K\mathcal{H}\!om_A(E,F)\to F$ 定义的 $A$-\Term{linear}{线性}\Term{arrow}{箭头}
\[E\otimes_K\mathcal{H}\!om_A(E,F)\otimes_A G\longrightarrow F\otimes_A G.\]
\Term{arrow}{箭头} 5.4.1 \Term{derive}{导出}为 $D(K)$ 中一个称为\Term{tensor-product}{张量积}\Term{arrow}{箭头}的\Term{arrow}{箭头}
\begin{equation}R\mathcal{H}\!om_A(E,F)\overset L\otimes_A G\longrightarrow R\mathcal{H}\!om_A(E,F\overset L\otimes_A G),\tag{5.4.2}\end{equation}
\SourcePage{179}
其中 $E\in\operatorname{ob}D^-(A)$、$F\in\operatorname{ob}D^+(A^e)$、$G\in\operatorname{ob}D^-(A)$，且 $F\overset L\otimes_A G\in\operatorname{ob}D^+(A)$。事实上，若 $E,G$ \Term{bounded-above}{上有界}且各\Term{set-component}{分量}\Term{flat}{平坦}，$F$ \Term{bounded-below}{下有界}且各\Term{set-component}{分量}\Term{injective}{内射}，则 5.4.2 的第一项等同于 $\mathcal{H}\!om_A^\bullet(E,F)\otimes_A G$；由 5.4.1 与自然\Term{arrow}{箭头}
\[\mathcal{H}\!om_A^\bullet(E,F\otimes_A G)\longrightarrow R\mathcal{H}\!om_A(E,F\otimes_A G)=R\mathcal{H}\!om_A(E,F\overset L\otimes_A G)\]
的复合，将它映入第二项。

\Term{arrow}{箭头} 5.4.2 一般不是\Term{isomorphism}{同构}。不过在某些重要情形它是\Term{isomorphism}{同构}，例如 $E$ \Term{perfect-complex}{完美}时（归结到 $E=A$ 即可验证）。后面还将看到另一个有趣的例子（6.2.3）。

\Term{arrow}{箭头} 5.3.5 和 5.4.2 将在本讲后文中起关键\Term{action}{作用}。我们想使用的\Term{object}{对象} $F$ 的例子是：a) $F=A$，或者 $F=A$ 以自然方式视为右 $A$-\Term{module}{模}，并通过一个 $K$-\Term{algebra}{代数}\Term{homomorphism}{同态} $\varphi:A\to A$ 视为左 $A$-\Term{module}{模}（参见 5.13）；b) $F=M\overset L\otimes_K A$，其中 $M\in\operatorname{ob}D^+(K)$（第 6 节中，$M$ 将是一个\Term{dualizing-complex}{对偶化复形}，特别参见 6.3.2、6.4.1、6.5.6）。

\textbf{5.5.} 以 $DF(A)$ 表示带\Term{finite}{有限}\Term{filtration}{滤过}的左 $A$-\Term{module}{模}\Term{complex}{复形}的\Term{derived-category}{导出范畴}（[3] V）。在 5.2 的情形，对 $E\in\operatorname{ob}D^-F(A)$、$F\in\operatorname{ob}D^-F(B,A)\ (=D^-F(B\otimes_K A^\circ))$、$G\in\operatorname{ob}D^+F(B)$，有 $DF(K)$ 中类似于 5.2.2 的\Term{canonical}{典范}\Term{isomorphism}{同构}。同样有类似于 5.3.4、5.3.5、5.3.6、5.4.2 的“\Term{filtered}{滤过}\Term{derive}{导出}”\Term{arrow}{箭头}。这些\Term{arrow}{箭头}与取\Term{associated-graded}{相伴分次对象}相容，含义如 [3] V 1、2。

\textbf{5.6.} 假设 $A$ 在 $K$ 上\Term{flat}{平坦}。以 $D(A)_{\mathrm{parf}}$ 表示 $D(A)$ 的\Term{full-subcategory}{满子范畴}，其\Term{object}{对象}是\Term{finite}{有限}\Term{perfect-amplitude}{完美振幅}的\Term{complex}{复形}，即\Term{perfect-complex}{完美}且有\Term{finite-tor-dimension}{有限 Tor 维数}的\Term{complex}{复形}（SGA 6 I）。对 $E\in\operatorname{ob}D(A)_{\mathrm{parf}}$、$F\in\operatorname{ob}D^b(A^e)$，把 $G=E$ 时的 5.4.2 的逆\Term{isomorphism}{同构}与\Term{evaluation-map}{求值箭头} 5.3.5 复合，定义称为\Term{trace}{迹}的\Term{canonical}{典范}\Term{arrow}{箭头}
\begin{equation}\operatorname{Tr}_A:R\mathcal{H}\!om_A(E,F\overset L\otimes_A E)\longrightarrow F\overset L\otimes_{A^e}A.\tag{5.6.1}\end{equation}
由此得到\Term{trace}{迹}\Term{arrow}{箭头}
\SourcePage{180}
\begin{equation}\operatorname{Tr}_A:\operatorname{Hom}_A(E,F\overset L\otimes_A E)\longrightarrow H^0(T,F\overset L\otimes_{A^e}A).\tag{5.6.2}\end{equation}
当 $F=A$ 时，再与\Term{canonical}{典范}\Term{arrow}{箭头} $A\overset L\otimes_{A^e}A\to A\otimes_{A^e}A=A_\natural$（5.1）复合，得到仍称为\Term{trace}{迹}的\Term{arrow}{箭头}
\begin{equation}\operatorname{Tr}_A:R\mathcal{H}\!om_A(E,E)\longrightarrow A_\natural,\tag{5.6.3}\end{equation}
\begin{equation}\operatorname{Tr}_A:\operatorname{Hom}_A(E,E)\longrightarrow H^0(T,A_\natural).\tag{5.6.4}\end{equation}
注意，5.6.4 是一个相当平常、具有\Term{local-nature}{局部性质}的\Term{arrow}{箭头}，而 5.6.3 则是一个“实质性”的、具有整体性质的\Term{arrow}{箭头}。

以 $DF(A)_{\mathrm{parf}}$ 表示 $DF(A)$ 的\Term{full-subcategory}{满子范畴}，其\Term{object}{对象}是\Term{graded-object}{分次对象}具有\Term{finite}{有限}\Term{perfect-amplitude}{完美振幅}的\Term{filtered-complex}{滤过复形}。于是，对 $E\in\operatorname{ob}DF(A)_{\mathrm{parf}}$、$F\in\operatorname{ob}D^bF(A^e)$，有 $DF(K)$ 中类似于 5.6.1 的\Term{arrow}{箭头}，从而有类似于 5.6.2、5.6.3、5.6.4 的\Term{arrow}{箭头}。由 5.5 中的说明，\Term{filtered}{滤过}版本的\Term{arrow}{箭头} 5.6.1 与取\Term{associated-graded}{相伴分次对象}相容（参见 [3] V 3.7.2）；与该处相同地\Term{pushout}{推出}，若 $F$ 的\Term{filtration}{滤过}集中于\Term{degree}{次数} $0$，则对 $u\in\operatorname{Hom}_{DF(A)}(E,F\overset L\otimes_A E)$ 有
\begin{equation}\operatorname{Tr}_A(u)=\sum_i\operatorname{Tr}_A\operatorname{gr}^i u.\tag{5.6.5}\end{equation}
这个公式特别蕴含，对 $E\in\operatorname{ob}D(A)_{\mathrm{parf}}$、$F\in\operatorname{ob}D^b(A^e)$、$u\in\operatorname{Hom}_A(E,F\overset L\otimes_A E)$ 有
\begin{equation}\operatorname{Tr}(u[1])=-\operatorname{Tr}(u),\tag{5.6.6}\end{equation}
其中 $u[1]:E[1]\to(F\overset L\otimes_A E)[1]$ 是 $u$ 的\Term{translation}{平移}。

\textbf{5.7.} 5.6 中定义的\Term{trace}{迹}同时推广了“通常”的交换\Term{trace}{迹}（SGA 6 I；[3] V 3）和 Stallings [5] 及 Grothendieck 的“\Term{noncommutative-trace}{非交换迹}”。

a) 假设 $A=K$。这时 $A^e=A_\natural=K$，\Term{arrow}{箭头} 5.6.1 写成
\[\operatorname{Tr}_K:R\mathcal{H}\!om_K(E,F\overset L\otimes_K E)\longrightarrow F\]
\SourcePage{181}
并与 [3] V 3.7.1 中定义的\Term{trace}{迹}\Term{arrow}{箭头}一致（这里是该处的\Term{complex}{复形} $K$ 化为 $A$ 的特殊情形）。特别地，$\operatorname{Tr}_K:R\mathcal{H}\!om_K(E,E)\to K$ 与 SGA 6 I 8 中定义的\Term{trace}{迹}一致。

b) 再假设 $A$ 是一个\Term{flat}{平坦} $K$-\Term{algebra}{代数}（未必交换）。我们要对严格\Term{perfect-complex}{完美}的 $E$ 明确写出 5.6.4（即 $E$ \Term{bounded}{有界}，各\Term{set-component}{分量}\Term{local-scheme}{局部}为\Term{finite-type}{有限型}\Term{free-module}{自由模}的\Term{direct-summand}{直和因子}）。先取 $E=A$；通过 $f\mapsto f(1)$、$a\mapsto(x\mapsto xa)$，$R\mathcal{H}\!om_A(A,A)=\mathcal{H}\!om_A(A,A)$ 等同于 $A$，而 5.6.3 等同于\Term{canonical}{典范}\Term{projection}{投影} $A\to A_\natural$，记为 $a\mapsto a_\natural$。所以若 $d_a$ 表示右乘 $a$，则
\begin{equation}\operatorname{Tr}_A(d_a)=a_\natural.\tag{5.7.1}\end{equation}
对 $1\leq i\leq n$，设 $E_i$ 是 $A$ 的一个副本，并令 $E=\bigoplus_{1\leq i\leq n}E_i$。$F=A$、$G=E$ 时的\Term{isomorphism}{同构} 5.4.2 写成
\[\bigoplus_{1\leq i,j\leq n}\mathcal{H}\!om_A(E_i,A)\otimes_A E_j\longrightarrow\mathcal{H}\!om_A(E,E).\]
由\Term{evaluation-map}{求值箭头} 5.3.5 的定义，对 $u=\sum u_{ij}\otimes1\in\mathcal{H}\!om_A(E,E)$（$u_{ij}\otimes1\in\mathcal{H}\!om_A(E_i,A)\otimes_A E_j$）有
\begin{equation}\operatorname{Tr}_A(u)=\sum_i(u_{ii})_\natural.\tag{5.7.2}\end{equation}
现在取 $E$ 为 $A^n$ 的一个\Term{direct-summand}{直和因子}，是某个\Term{projector}{投影算子} $p$ 的像，以 $i:E\to A^n$ 表示包含。则对 $u\in\mathcal{H}\!om_A(E,E)$ 有
\begin{equation}\operatorname{Tr}_A(u)=\operatorname{Tr}_A(iup),\tag{5.7.3}\end{equation}
这既可由 5.6.5 看出，也可直接从定义看出。最后，假设 $E$ 严格\Term{perfect-complex}{完美}；若 $u\in\operatorname{Hom}_A(E,E)$ 是\Term{complex}{复形}的\Term{homomorphism}{同态}，由 5.6.5、5.6.6 得
\begin{equation}\operatorname{Tr}_A(u)=\sum_i(-1)^i\operatorname{Tr}_A(u^i).\tag{5.7.4}\end{equation}
\SourcePage{182}
5.6.4 是对一个\Term{complex}{复形}\Term{homomorphism}{同态} $\mathcal{H}\!om_A^\bullet(E,E)\to A$（5.6.3）应用 $H^0$ 而得到的；这一事实蕴含，若 $u,v\in\operatorname{Hom}_A(E,E)$ \Term{homotopic}{同伦}，则
\begin{equation}\operatorname{Tr}_A(u)=\operatorname{Tr}_A(v).\tag{5.7.5}\end{equation}
由于 $D(A)$ 中 $E$ 的每个\Term{endomorphism}{自同态}都\Term{local-scheme}{局部}由一个\Term{complex}{复形}\Term{endomorphism}{自同态}定义，且\Term{local-scheme}{局部}在\Term{homotopy}{同伦}意义下唯一（SGA 6 I 4.10），公式 5.7.2 至 5.7.4 因而可用于计算 $D(A)$ 中 $E$ 的任意\Term{endomorphism}{自同态} $u$ 的 $\operatorname{Tr}_A(u)$；与 5.7.5 结合，就给出了 $\operatorname{Tr}_A(u)$ 的直接定义。当\Term{topos}{拓扑斯} $T$ 为点\Term{topos}{拓扑斯}时，便得到 SGA $4\frac12$ Rapport 4.3 的定义。

\textbf{5.8.} 设 $P,Q$ 是 $A$-\Term{bimodule}{双模}。则 $P\otimes_A Q$ 和 $Q\otimes_A P$ 也是 $A$-\Term{bimodule}{双模}，并有\Term{canonical}{典范}\Term{isomorphism}{同构}
\begin{equation}(P\otimes_A Q)\otimes_{A^e}A\xrightarrow\sim P\otimes_{A^e}Q\xrightarrow\sim(Q\otimes_A P)\otimes_{A^e}A\xrightarrow\sim A\otimes_{A^e}(P\otimes_K Q)\otimes_{A^e}A.\tag{5.8.1}\end{equation}
上述各\Term{object}{对象}都\Term{canonical}{典范}地等同于 $P\otimes_K Q$ 对由形如 $axb\otimes y-x\otimes bya$ 的\Term{section}{截面}生成的\Term{submodule}{子模}的商 $K$-\Term{module}{模}，其中 $y\in P$、$y\in Q$、$a,b\in K$。当 $A$ 在 $K$ 上\Term{flat}{平坦}时，这些\Term{isomorphism}{同构}\Term{derive}{导出}为\Term{isomorphism}{同构}
\begin{equation}(P\overset L\otimes_A Q)\overset L\otimes_{A^e}A\xrightarrow\sim P\overset L\otimes_{A^e}Q\xrightarrow\sim(Q\overset L\otimes_A P)\overset L\otimes_{A^e}A\xrightarrow\sim A\overset L\otimes_{A^e}(P\overset L\otimes_K Q)\overset L\otimes_{A^e}A,\tag{5.8.2}\end{equation}
其中 $P,Q\in\operatorname{ob}D^-(A^e)$。另一方面，\Term{module}{模}层次的复合\Term{arrow}{箭头}
\[\begin{aligned}
\mathcal{H}\!om_A(E,P\otimes_A F)\otimes_K\mathcal{H}\!om_A(F,Q\otimes_A E)&\longrightarrow\mathcal{H}\!om_A(E,P\otimes_A Q\otimes_A E),\\
f\otimes g&\longmapsto(P\otimes_A g)f
\end{aligned}\]
\Term{derive}{导出}为
\begin{equation}R\mathcal{H}\!om_A(E,P\overset L\otimes_A F)\overset L\otimes_K R\mathcal{H}\!om_A(F,Q\overset L\otimes_A E)\longrightarrow R\mathcal{H}\!om_A(E,P\overset L\otimes_A Q\overset L\otimes_A E),\tag{5.8.3}\end{equation}
\SourcePage{183}
其中 $P,Q\in\operatorname{ob}D^b(A^e)$、$E,F\in\operatorname{ob}D^b(A)$，并要求上式出现的\Term{tensor-product}{张量积}及 $R\mathcal{H}\!om$ 均在 $D^b$ 中。当 $E,F$ \Term{perfect-complex}{完美}时，通过 5.4.2，这个\Term{arrow}{箭头}等同于从\Term{evaluation-map}{求值箭头} $F\overset L\otimes_K R\mathcal{H}\!om_A(F,Q)\to Q$ \Term{derive}{导出}的\Term{arrow}{箭头}。容易由此\Term{pushout}{推出}有\Term{commutative-diagram}{交换方块}
\begin{equation}
\resizebox{.94\linewidth}{!}{\begin{tikzcd}[ampersand replacement=\&,column sep=large,row sep=large]
R\mathcal{H}\!om_A(E,P\overset L\otimes_A F)\overset L\otimes_K R\mathcal{H}\!om_A(F,Q\overset L\otimes_A E) \arrow[r,"(1)"]\arrow[d,"(2)"'] \& R\mathcal{H}\!om_A(E,P\overset L\otimes_A Q\overset L\otimes_A E)\arrow[d,"\operatorname{Tr}_A"]\\
R\mathcal{H}\!om_A(F,Q\overset L\otimes_A P\overset L\otimes_A F)\arrow[r,"\operatorname{Tr}_A"'] \& P\overset L\otimes_{A^e}Q
\end{tikzcd}}
\tag{5.8.4}\end{equation}
其中 (1)、(2) 由 5.8.3 给出（关键在于，5.8.4 中的两个复合\Term{arrow}{箭头}都由\Term{evaluation-map}{求值箭头} $F\overset L\otimes_K R\mathcal{H}\!om_A(F,Q)\to Q$、$E\overset L\otimes_K R\mathcal{H}\!om_A(E,P)\to P$ 的\Term{tensor-product}{张量积}定义）。因此，对 $u\in\operatorname{Hom}_A(E,P\overset L\otimes_A F)$、$v\in\operatorname{Hom}(F,Q\overset L\otimes_A E)$ 有
\begin{equation}\operatorname{Tr}_A(vu)=\operatorname{Tr}_A(uv).\tag{5.8.5}\end{equation}
有时以 $\langle u,v\rangle_A$ 表示 5.8.5 两边的共同值。

\SGAsection{5.9}{迹与标量扩张}
设 $A\to B$ 是\Term{flat}{平坦} $K$-\Term{algebra}{代数}的\Term{morphism}{态射}。由此得到一个 $K$-\Term{algebra}{代数}\Term{morphism}{态射} $A^e\to B^e$，从而有\Term{extension-of-scalars}{标量扩张}\Term{functor}{函子} $-\overset L\otimes_{A^e}B^e:D(A^e)\to D(B^e)$。对 5.3.5 中的 $E,F$，若 $F\overset L\otimes_{A^e}B^e\in\operatorname{ob}D^b(B^e)$ 且 $R\mathcal{H}\!om_B(B\overset L\otimes_A E,F\overset L\otimes_{A^e}B^e)\in\operatorname{ob}D^b(B^\circ)$，则有 $D(K)$ 中的\Term{commutative-diagram}{交换方块}
\begin{equation}
\resizebox{.94\linewidth}{!}{\begin{tikzcd}[ampersand replacement=\&,column sep=large,row sep=large]
R\mathcal{H}\!om_A(E,F)\overset L\otimes_A E\arrow[r,"(1)"]\arrow[d,"(3)"'] \& F\overset L\otimes_{A^e}A\arrow[d,"(4)"]\\
R\mathcal{H}\!om_B(B\overset L\otimes_A E,F\overset L\otimes_{A^e}B^e)\overset L\otimes_B(B\overset L\otimes_A E)\arrow[r,"(2)"'] \& (F\overset L\otimes_{A^e}B^e)\overset L\otimes_{B^e}B
\end{tikzcd}}
\tag{5.9.1}\end{equation}
其中\Term{arrow}{箭头} (1)、(2) 由 5.3.6 定义，(4) 由 $A\to B$ 定义，(3) 由\Term{canonical}{典范}\Term{arrow}{箭头} $F\to F\overset L\otimes_{A^e}B^e$ 定义。验证是直接的。同样，对 5.4.2 中的 $E,F,G$，若 $F\overset L\otimes_{A^e}B^e\in\operatorname{ob}D^+(B^e)$、$(F\overset L\otimes_{A^e}B^e)\overset L\otimes_A G\in\operatorname{ob}D^+(B)$，则有
\SourcePage{184}
$D(K)$ 中的\Term{commutative-diagram}{交换方块}
\begin{equation}
\resizebox{.94\linewidth}{!}{\begin{tikzcd}[ampersand replacement=\&,column sep=large,row sep=large]
R\mathcal{H}\!om_A(E,F)\overset L\otimes_A G\arrow[r,"(1)"]\arrow[d,"(3)"'] \& R\mathcal{H}\!om_A(E,F\overset L\otimes_A G)\arrow[d,"(4)"]\\
R\mathcal{H}\!om_B(B\overset L\otimes_A E,F\overset L\otimes_{A^e}B^e)\overset L\otimes_B(B\overset L\otimes_A G)\arrow[r,"(2)"'] \& R\mathcal{H}\!om_B(B\overset L\otimes_A E,(F\overset L\otimes_{A^e}B^e)\overset L\otimes_B(B\overset L\otimes_A G))
\end{tikzcd}}
\tag{5.9.2}\end{equation}
其中 (1)、(2) 由 5.4.2 定义，(3)、(4) 由 $F\to F\overset L\otimes_{A^e}B^e$ 定义。因此，对 5.6.1 中的 $E,F$，若 $F\overset L\otimes_{A^e}B^e\in\operatorname{ob}D^b(B^e)$，则有一个与前面类似的 $D(K)$ 中的\Term{commutative-diagram}{交换方块}，其竖直\Term{arrow}{箭头}由 5.6.1 定义：
\begin{equation}
\resizebox{.94\linewidth}{!}{\begin{tikzcd}[ampersand replacement=\&,column sep=large,row sep=large]
R\mathcal{H}\!om_A(E,F\overset L\otimes_A E)\arrow[r,"\operatorname{Tr}_A"]\arrow[d] \& F\overset L\otimes_{A^e}A\arrow[d]\\
R\mathcal{H}\!om_B(B\overset L\otimes_A E,(F\overset L\otimes_{A^e}B^e)\overset L\otimes_B(B\overset L\otimes_A E))\arrow[r,"\operatorname{Tr}_B"'] \& (F\overset L\otimes_{A^e}B^e)\overset L\otimes_{B^e}B
\end{tikzcd}}
\tag{5.9.3}\end{equation}
最后，若给定 $F'\in\operatorname{ob}D^b(B^e)$ 使 $R\mathcal{H}\!om_A(E,F')\in\operatorname{ob}D^b(B^\circ)$，以及 $f:F\overset L\otimes_{A^e}B^e\to F'$，则有比 5.9.1 略为一般的相容性：它由一个类似的\Term{commutative-diagram}{交换方块}表达，在下行以 $F'$ 代替 $F\overset L\otimes_{A^e}B^e$，并由 $f$ 定义竖直\Term{arrow}{箭头}。同样，以 $F'$ 代替 $F\overset L\otimes_{A^e}B^e$，可得到类似于 5.9.2 和 5.9.3 的\Term{commutative-diagram}{交换方块}。特别取 $F=A$，并取 $A\to B$ 所定义的\Term{arrow}{箭头} $f:A\overset L\otimes_{A^e}B^e\to B$，便有\Term{commutative-diagram}{交换方块}
\begin{equation}\begin{tikzcd}[column sep=large,row sep=large]
R\mathcal{H}\!om_A(E,E)\arrow[r,"\operatorname{Tr}_A"]\arrow[d] & A\overset L\otimes_{A^e}A\arrow[d]\\
R\mathcal{H}\!om_B(B\overset L\otimes_A E,B\overset L\otimes_A E)\arrow[r,"\operatorname{Tr}_B"'] & B\overset L\otimes_{B^e}B,
\end{tikzcd}\tag{5.9.4}\end{equation}
其中左（相应地右）竖直\Term{arrow}{箭头}由\Term{extension-of-scalars}{标量扩张}给出（相应地为从 $A\to B$ 得到的\Term{canonical}{典范}\Term{arrow}{箭头}）。

\SourcePage{185}
\SGAsection{5.10}{迹与标量限制}
设 $A\to B$ 是\Term{flat}{平坦} $K$-\Term{algebra}{代数}的\Term{morphism}{态射}，且 $B$ 作为左 $A$-\Term{module}{模}\Term{flat}{平坦}而\Term{finite-presentation}{有限表现}（即\Term{local-scheme}{局部}为\Term{finite-type}{有限型}自由 $A$-\Term{module}{模}的\Term{direct-summand}{直和因子}）。例如，$A\to B$ 是由\Term{finite-index}{有限指数}的\Term{discrete}{离散}群单\Term{homomorphism}{同态} $H\hookrightarrow G$ 定义的\Term{morphism}{态射} $K[H]\to K[G]$ 时就是如此。\Term{restriction-of-scalars}{标量限制}\Term{homomorphism}{同态} $\mathcal{H}\!om_B(B,B)\to\mathcal{H}\!om_A(B,B)$（这里将 $B$ 视为左\Term{module}{模}），通过\Term{canonical}{典范}\Term{isomorphism}{同构} $B\xrightarrow\sim\mathcal{H}\!om_B(B,B)$、$\mathcal{H}\!om_A(B,A)\otimes_A B\xrightarrow\sim\mathcal{H}\!om_A(B,B)$，等同于一个（$B$-\Term{bimodule}{双模}的）\Term{arrow}{箭头}
\begin{equation}B\longrightarrow\mathcal{H}\!om_A(B,A)\otimes_A B.\tag{5.10.1}\end{equation}
5.10.1 与\Term{evaluation-map}{求值箭头} $\mathcal{H}\!om_A(B,A)\otimes_A B\to A_\natural$（5.3.3）的复合，就是\Term{homomorphism}{同态} $f\mapsto\operatorname{Tr}_A(f)$（定义为 $\operatorname{Tr}_A(d_f)$，其中 $d_f$ 是右乘 $f$）。因此，由公式 $\operatorname{Tr}_A(fg)=\operatorname{Tr}_A(gf)$（5.8.5），它通过 $B_\natural$ \Term{resolution}{分解}为一个仍记为 $\operatorname{Tr}_{B/A}:B_\natural\to A_\natural$ 的\Term{arrow}{箭头}。所以有\Term{commutative-diagram}{交换方块}
\begin{equation}\begin{tikzcd}[column sep=large,row sep=large]
B\arrow[r,"\operatorname{Tr}_B"]\arrow[d,"5.10.1"'] & B_\natural\arrow[d,"\operatorname{Tr}_{B/A}"]\\
\mathcal{H}\!om_A(B,A)\otimes_A B\arrow[r,"\mathrm{ev}"'] & A_\natural.
\end{tikzcd}\tag{5.10.2}\end{equation}
在 $({}_BE,{}_BP_A)$ 的情形，有右 $B$-\Term{module}{模}的\Term{homomorphism}{同态}
\begin{equation}\mathcal{H}\!om_B(E,P)\otimes_A B\longrightarrow\mathcal{H}\!om_B(E,P\otimes_A B),\qquad u\otimes b\longmapsto(x\longmapsto u(x)\otimes b),\tag{5.10.3}\end{equation}
它是\Term{isomorphism}{同构}，这一点立即可证（只须看其底下的 $K$-\Term{linear}{线性}\Term{homomorphism}{同态}是\Term{isomorphism}{同构}，而这由 $B$ \Term{local-scheme}{局部}为\Term{finite-type}{有限型}自由 $A$-\Term{module}{模}的\Term{direct-summand}{直和因子}而得）。此\Term{isomorphism}{同构}\Term{derive}{导出}为 $D(B)$ 中的\Term{isomorphism}{同构}
\begin{equation}R\mathcal{H}\!om_B(E,P)\otimes_A B\xrightarrow\sim R\mathcal{H}\!om_B(E,P\otimes_A B),\tag{5.10.4}\end{equation}
\SourcePage{186}
其中 $E\in\operatorname{ob}D^-(B)$、$P\in\operatorname{ob}D^+(B\otimes_K A^\circ)$（以形如 $B\otimes_A L$ 的\Term{module}{模}\Term{resolution}{消解} $E$，其中 $L$ \Term{flat}{平坦}；以\Term{injective}{内射} $(B\otimes_K A^\circ)$-\Term{module}{模}\Term{resolution}{消解} $P$，由于 $A$ 在 $K$ 上\Term{flat}{平坦}，后者作为 $B$-\Term{module}{模}也必然\Term{injective}{内射}：这使我们能在左边“去掉 $R$”，而对 $E$ 所取的\Term{resolution}{消解}也使我们能在右边去掉它）。

对 $E\in\operatorname{ob}D^-(B)$、$M\in\operatorname{ob}D^+(K)$、$F\in\operatorname{ob}D^-(B)$，定义 $D(X)$ 中的一个\Term{arrow}{箭头}\footnote{在第 6 节中，将把以下构造应用于 $M$ 是\Term{dualizing-complex}{对偶化复形}的情形（6.7）。}
\begin{equation}R\mathcal{H}\!om_B(E,B\overset L\otimes_K M)\overset L\otimes_B F\longrightarrow R\mathcal{H}\!om_A(E,A\overset L\otimes_K M)\overset L\otimes_A F,\tag{5.10.5}\end{equation}
它是以下\Term{arrow}{箭头}的复合：
\[\begin{aligned}
{}&R\mathcal{H}\!om_B(E,B\overset L\otimes_K M)\overset L\otimes_B F\\
&\quad\xrightarrow{(1)}R\mathcal{H}\!om_B(E,\mathcal{H}\!om_A(B,A)\otimes_A B\overset L\otimes_K M)\overset L\otimes_B F\\
&\quad\xrightarrow{(2)}R\mathcal{H}\!om_B(E,\mathcal{H}\!om_A(B,A)\overset L\otimes_K M)\otimes_A B\overset L\otimes_B F\\
&\quad\xrightarrow{(3)}R\mathcal{H}\!om_B(E,\mathcal{H}\!om_A(B,A\overset L\otimes_K M))\overset L\otimes_A F\\
&\quad\xrightarrow{(4)}R\mathcal{H}\!om_A(E,A\overset L\otimes_K M)\overset L\otimes_A F,
\end{aligned}\]
其中 (1) 由 5.10.1 定义，(2) 由 $P=\mathcal{H}\!om_A(B,A)\overset L\otimes_K M$ 时的 5.10.4 的逆\Term{isomorphism}{同构}定义，(3) 由\Term{canonical}{典范}\Term{isomorphism}{同构} $\mathcal{H}\!om_A(B,A)\overset L\otimes_K M\xrightarrow\sim\mathcal{H}\!om_A(B,A\overset L\otimes_K M)$（5.4.2）定义，(4) 由\Term{adjunction}{伴随}\Term{isomorphism}{同构} 5.2.2 定义。

\begin{SGAresult}{命题}{5.10.6}
假设 $M\overset L\otimes_K F\in\operatorname{ob}D^+(B)$。则有 $D(K)$ 中的\Term{commutative-diagram}{交换方块}
\begin{equation}\begin{tikzcd}[column sep=large,row sep=large]
R\mathcal{H}\!om_B(E,B\overset L\otimes_K M)\overset L\otimes_B F\arrow[r,"5.4.2"]\arrow[d,"5.10.5"'] & R\mathcal{H}\!om_B(E,M\overset L\otimes_K F)\arrow[d]\\
R\mathcal{H}\!om_A(E,A\overset L\otimes_K M)\overset L\otimes_A F\arrow[r,"5.4.2"'] & R\mathcal{H}\!om_A(E,M\overset L\otimes_K F),
\end{tikzcd}\tag{5.10.7}\end{equation}
其中右侧竖直\Term{arrow}{箭头}是\Term{restriction-of-scalars}{标量限制}。
\end{SGAresult}
\begin{proof}
\SourcePage{187}
需要验证以下方块\Term{commutative-diagram}{交换}：
\[\begin{tikzcd}[column sep=large,row sep=large]
E\overset L\otimes_K R\mathcal{H}\!om_B(E,B\overset L\otimes_K M)\overset L\otimes_B F\arrow[r]\arrow[d,"E\overset L\otimes 5.10.5"'] & M\overset L\otimes_K F\arrow[d,"\text{标量限制}"]\\
E\overset L\otimes_K R\mathcal{H}\!om_A(E,A\overset L\otimes_K M)\overset L\otimes_A F\arrow[r] & M\overset L\otimes_K F,
\end{tikzcd}\]
其中水平\Term{arrow}{箭头}由求值（5.3.6）给出。该方块是将 $-\overset L\otimes_B F$ 应用于 $F=B$ 时的类似方块而得，所以可假设 $F=B$。在此情形，5.10.7 的\Term{commutative-diagram}{交换性}\Term{equivalence}{等价}于三角形
\[\begin{tikzcd}[column sep=large,row sep=large]
R\mathcal{H}\!om_B(E,B\overset L\otimes_K M)\arrow[d,"(1)"']\arrow[dr,"(2)"] & \\
R\mathcal{H}\!om_B(E,\mathcal{H}\!om_A(B,B\overset L\otimes_K M))\arrow[r,"(3)"'] & R\mathcal{H}\!om_A(E,B\overset L\otimes_K M)
\end{tikzcd}\qquad(*)\]
的\Term{commutative-diagram}{交换性}。其中 (1) 由 5.10.1 定义，(2) 是\Term{restriction-of-scalars}{标量限制}，(3) 是\Term{adjunction}{伴随}\Term{isomorphism}{同构}。而由 5.10.1 给出的\Term{arrow}{箭头} $B\overset L\otimes_K M\to\mathcal{H}\!om_A(B,B\overset L\otimes_K M)$ 正是对应于\Term{restriction-of-scalars}{标量限制}\Term{functor}{函子}与 $\mathcal{H}\!om_A(B,-)$ 的\Term{adjunction}{伴随}\Term{arrow}{箭头}，故立即得到 $(*)$ 的\Term{commutative-diagram}{交换性}。
\end{proof}

设 $E\in\operatorname{ob}D^-(B)$、$M\in\operatorname{ob}D^b(K)$，使 $R\mathcal{H}\!om_B(E,B\overset L\otimes_K M)\in\operatorname{ob}D^b(B^\circ)$。由 5.3.5 有\Term{evaluation-map}{求值箭头}
\[R\mathcal{H}\!om_B(E,B\overset L\otimes_K M)\overset L\otimes_B E\longrightarrow M\overset L\otimes_K(B\overset L\otimes_{B^e}B),\]
再与\Term{canonical}{典范}\Term{arrow}{箭头} $B\overset L\otimes_{B^e}B\to B_\natural$ 复合，就得到\Term{arrow}{箭头}
\begin{equation}R\mathcal{H}\!om_B(E,B\overset L\otimes_K M)\overset L\otimes_B E\longrightarrow M\overset L\otimes_K B_\natural.\tag{5.10.8}\end{equation}

\begin{SGAresult}{命题}{5.10.9}
假设 $A$ 作为 $K$-\Term{module}{模}\Term{finite-presentation}{有限表现}且\Term{flat}{平坦}，$A_\natural,B_\natural$ 在 $K$ 上\Term{flat}{平坦}。设 $E\in\operatorname{ob}D^-(B)$、$M\in\operatorname{ob}D^b(K)$，且 $R\mathcal{H}\!om_B(E,B\overset L\otimes_K M)\in\operatorname{ob}D^b(B^\circ)$。则 $R\mathcal{H}\!om_A(E,A\overset L\otimes_K M)\in\operatorname{ob}D^b(A^\circ)$，并有\Term{commutative-diagram}{交换方块}
\SourcePage{188}
\begin{equation}\begin{tikzcd}[column sep=large,row sep=large]
R\mathcal{H}\!om_B(E,B\overset L\otimes_K M)\overset L\otimes_B E\arrow[r,"5.10.8"]\arrow[d,"5.10.5"'] & M\otimes_K B_\natural\arrow[d,"M\otimes_K\operatorname{Tr}_{B/A}"]\\
R\mathcal{H}\!om_A(E,A\overset L\otimes_K M)\overset L\otimes_A E\arrow[r,"5.10.8"'] & M\otimes_K A_\natural.
\end{tikzcd}\tag{5.10.10}\end{equation}
\end{SGAresult}
（对 $A,B$ 的附加假设——例如当 $A,B$ 为\Term{finite-group}{有限群}的 $K$-\Term{group-algebra}{群代数}时成立——大概是不必要的；应当有一个类似于 5.10.10 的\Term{commutative-diagram}{交换方块}，其中以 $A\overset L\otimes_{A^e}A$（相应地 $B\overset L\otimes_{B^e}B$）代替 $A_\natural$（相应地 $B_\natural$）。）
\begin{proof}
5.10.9 的第一个断言是显然的。证明第二个断言。可假定 $M$ 的各\Term{set-component}{分量}\Term{injective}{内射}，$E$ 的各\Term{set-component}{分量}形如 $B\otimes_K L$，其中 $L$ \Term{flat}{平坦}。于是有
\[\begin{aligned}
R\mathcal{H}\!om_B(E,B\overset L\otimes_K M)\overset L\otimes_B E&=\mathcal{H}\!om_B^\bullet(E,B\otimes_K M)\otimes_B E,\\
R\mathcal{H}\!om_A(E,A\overset L\otimes_K M)\overset L\otimes_A E&=\mathcal{H}\!om_A^\bullet(E,A\otimes_K M)\otimes_A E,
\end{aligned}\]
而 5.10.10 的\Term{arrow}{箭头}均定义在\Term{complex}{复形}层次。因此归结为验证：若 $L,M$ 是 $K$-\Term{module}{模}，$E$ 为左 $B$-\Term{module}{模}，则有\Term{commutative-diagram}{交换方块}
\[\resizebox{.94\linewidth}{!}{\begin{tikzcd}[ampersand replacement=\&,column sep=large,row sep=large]
\mathcal{H}\!om_B(B\otimes_K L,B\otimes_K M)\otimes_B(B\otimes_K L)\arrow[r,"5.3.3"]\arrow[d,"(1)"'] \& M\otimes_K B_\natural\arrow[d,"M\otimes_K\operatorname{Tr}_{B/A}"]\\
\mathcal{H}\!om_A(B\otimes_K L,A\otimes_K M)\otimes_A(B\otimes_K L)\arrow[r,"5.3.3"'] \& M\otimes_K A_\natural,
\end{tikzcd}}\quad(*)\]
其中\Term{arrow}{箭头} (1) 如下定义：由 5.10.1 将 $B\otimes_K M$ 映入 $\mathcal{H}\!om_A(B,A)\otimes_A B\otimes_K M$，再与一个 5.10.3 型的\Term{isomorphism}{同构}及一个\Term{adjunction}{伴随}\Term{isomorphism}{同构}复合。由于 $A$ 作为 $K$-\Term{module}{模}\Term{flat}{平坦}而\Term{finite-presentation}{有限表现}，“\Term{tensor-product}{张量积}”\Term{arrow}{箭头}
\begin{equation}B\otimes_K\mathcal{H}\!om_K(L,M)=\mathcal{H}\!om_B(B,B)\otimes_K\mathcal{H}\!om_K(L,M)\longrightarrow\mathcal{H}\!om_B(B\otimes_K L,B\otimes_K M)\tag{2}\end{equation}
和
\begin{equation}\mathcal{H}\!om_A(B,A)\otimes_K\mathcal{H}\!om_K(L,M)\longrightarrow\mathcal{H}\!om_A(B\otimes_K L,A\otimes_K M)\tag{3}\end{equation}
均为\Term{isomorphism}{同构}；通过它们，(1) 等同于 5.10.1 与 $\mathcal{H}\!om_K(L,M)\otimes_K L$ 的（在 $K$ 上的）\Term{tensor-product}{张量积}。另一方面，通过 (2)（相应地 (3)），$(*)$ 中上方（相应地下方）水平\Term{arrow}{箭头}等同于
\SourcePage{189}
$\operatorname{Tr}_B:B\to B_\natural$（相应地 $\mathrm{ev}:\mathcal{H}\!om_A(B,A)\otimes_A B\to A_\natural$）与\Term{evaluation-map}{求值箭头} $\mathcal{H}\!om_K(L,M)\otimes_K L\to M$ 的（在 $K$ 上的）\Term{tensor-product}{张量积}。所以 $(*)$ 的\Term{commutative-diagram}{交换性}由 5.10.2 的\Term{commutative-diagram}{交换性}\Term{pushout}{推出}，5.10.9 得证。
\end{proof}

\begin{SGAresult}{推论}{5.10.11}
在 5.10.9 的假设下，对 $E\in\operatorname{ob}D(B)_{\mathrm{parf}}$、$M\in\operatorname{ob}D^b(K)$，有 $D(K)$ 中的\Term{commutative-diagram}{交换方块}
\begin{equation}\begin{tikzcd}[column sep=large,row sep=large]
R\mathcal{H}\!om_B(E,M\overset L\otimes_K E)\arrow[r,"\operatorname{Tr}_B"]\arrow[d] & M\otimes_K B_\natural\arrow[d]\\
R\mathcal{H}\!om_A(E,M\overset L\otimes_K E)\arrow[r,"\operatorname{Tr}_A"'] & M\otimes_K A_\natural,
\end{tikzcd}\tag{5.10.12}\end{equation}
其中左侧竖直\Term{arrow}{箭头}是\Term{restriction-of-scalars}{标量限制}。
\end{SGAresult}
\begin{proof}事实上，只须结合 5.10.6 和 5.10.9。\end{proof}

\SGAsection{5.11}{有限群的应用：整除性质}
设 $G$ 为\Term{finite-group}{有限群}，\Term{identity-element}{单位元}记作 $e$，置 $K[G]=B$。$K$-\Term{module}{模} $B_\natural$ 是以 $G$ 的\Term{conjugacy-class}{共轭类}\Term{set}{集合} $G_\natural$ 为基的\Term{free-module}{自由模}；$B_\natural$ 的任意\Term{section}{截面} $s$ 唯一地写成 $s=\sum_{x\in G_\natural}s(x)x$。

\begin{SGAresult}{命题}{5.11.1}
（参见 SGA $4\frac12$ Rapport 4.5。）设 $E\in\operatorname{ob}D(B)_{\mathrm{parf}}$、$M\in\operatorname{ob}D^b(K)$。对 $u\in\operatorname{Hom}_B(E,M\overset L\otimes_K E)$，在 $H^0(T,M)$ 中有
\[\operatorname{Tr}_K(u)=|G|\operatorname{Tr}_B(u)(e).\footnotemark\]
\end{SGAresult}
\footnotetext{若 $X$ 是一个\Term{set}{集合}，以 $|X|$ 表示 $X$ 的\Term{cardinality}{基数}。}
\begin{proof}
事实上，只须取 $A=K$ 应用 5.10.11，并注意到 $\operatorname{Tr}_{B/K}:B_\natural\to K$ 由下式给出：
\begin{equation}\operatorname{Tr}_{B/K}(g)=\begin{cases}|G|&\text{若 }g=e,\\0&\text{若 }g\ne e.\end{cases}\tag{5.11.2}\end{equation}
\end{proof}

\begin{SGAresult}{命题}{5.11.3}
在 5.11.1 的假设下，设 $g\in G$，以 $Z_g$ 表示 $g$ 在 $G$ 中的\Term{centralizer}{中心化子}，置 $A=K[Z_g]$。设 $u\in\operatorname{Hom}_B(E,M\overset L\otimes_K E)$；以 $gu\in\operatorname{Hom}_A(E,M\overset L\otimes_K E)$ 表示 $u$ 在 $\operatorname{Hom}_A(E,M\overset L\otimes_K E)$ 中的像与左乘 $g$ 的复合。则在 $H^0(T,M)$ 中有
\SourcePage{190}
\[\operatorname{Tr}_A(gu)(e)=\operatorname{Tr}_B(u)(g^{-1}).\]
\end{SGAresult}
\begin{proof}
更一般地，我们将证明有\Term{commutative-diagram}{交换方块}
\begin{equation}\begin{tikzcd}[column sep=large,row sep=large]
R\mathcal{H}\!om_B(E,M\overset L\otimes_K E)\arrow[r,"\operatorname{Tr}_B"]\arrow[d] & M\otimes_K B_\natural\arrow[r,"x\mapsto x(g^{-1})"] & M\arrow[d,equal]\\
R\mathcal{H}\!om_A(E,M\overset L\otimes_K E)\arrow[r,"\operatorname{Tr}_A"'] & M\otimes_K A_\natural\arrow[r,"x\mapsto x(e)"'] & M,
\end{tikzcd}\tag{5.11.4}\end{equation}
其中左侧竖直\Term{arrow}{箭头}是\Term{restriction-of-scalars}{标量限制}与左乘 $g$ 的复合。该方块是以下图的外框：
\[\resizebox{.94\linewidth}{!}{\begin{tikzcd}[ampersand replacement=\&,column sep=large,row sep=large]
R\mathcal{H}\!om_B(E,M\overset L\otimes_K E)\arrow[r,"\operatorname{Tr}_B"]\arrow[d,"\mathrm{rest}"']\arrow[dr,phantom,"(1)" description] \& M\otimes_K B_\natural\arrow[r,"x\mapsto x(g^{-1})"]\arrow[d,"\operatorname{Tr}_{B/A}"] \& M\arrow[dd,equal]\\
R\mathcal{H}\!om_A(E,M\overset L\otimes_K E)\arrow[r,"\operatorname{Tr}_A"]\arrow[d,"{R\mathcal{H}\!om_A(E,s_g)}"']\arrow[dr,phantom,"(2)" description] \& M\otimes_K A_\natural\arrow[d,"s_g"]\arrow[dr,phantom,"(3)" description] \& \\
R\mathcal{H}\!om_A(E,M\overset L\otimes_K E)\arrow[r,"\operatorname{Tr}_A"'] \& M\otimes_K A_\natural\arrow[r,"x\mapsto x(e)"'] \& M,
\end{tikzcd}}\]
其中 $\mathrm{rest}$ 表示\Term{restriction-of-scalars}{标量限制}，$s_g$ 表示左乘 $g$。方块 (1) 由 5.10.11 \Term{commutative-diagram}{交换}，(2) 由 $\operatorname{Tr}_A$ 的定义\Term{commutative-diagram}{交换}。所以只须证明 (3) 的\Term{commutative-diagram}{交换性}；为此，只须证明对 $h\in G$ 有
\[\operatorname{Tr}_A(x\longmapsto gxh)(e)=h\text{ 在 }B_\natural\text{ 中的像在 }g^{-1}\text{ 处的值},\]
即
\begin{equation}\operatorname{Tr}_A(x\longmapsto gxh)(e)=\begin{cases}1&\text{若 }g^{-1}\text{ 与 }h\text{ 共轭},\\0&\text{否则}.\end{cases}\tag{*}\end{equation}
为此可将 $T$ 换成点\Term{topos}{拓扑斯}、将 $K$ 换成 $\mathbb Z$，此时只须证明 $(*)$ 乘以 $|Z_g|$ 后的等式。由 5.11.1，便归结为证明
\SourcePage{191}
\[\operatorname{Tr}_{\mathbb Z}(x\longmapsto gxh)=\begin{cases}|Z_g|&\text{若 }g^{-1}\text{ 与 }h\text{ 共轭},\\0&\text{否则}.\end{cases}\]
而这由左边是映射 $G\to G$、$x\mapsto gxh$ 的\Term{fixed-point}{不动点}个数而得。于是 5.11.4 \Term{commutative-diagram}{交换}，从而得到 5.11.3 的结论。
\end{proof}

\begin{SGAresult}{推论}{5.11.5}
在 5.11.3 的假设下，有
\[\operatorname{Tr}_K(gu)=|Z_g|\operatorname{Tr}_B(u)(g^{-1})=|Z_g|\operatorname{Tr}_A(gu)(e).\]
\end{SGAresult}

\SGAsection{5.12}{迹与外张量积}
对 $i=1,2$，设 $A_i$ 是一个\Term{flat}{平坦} $K$-\Term{algebra}{代数}，令 $A=A_1\otimes_K A_2$。设 $E_i\in\operatorname{ob}D^-(A_i)$、$F_i\in\operatorname{ob}D^b(A_i^e)$，使得 $R\mathcal{H}\!om_{A_i}(E_i,F_i)\in\operatorname{ob}D^b(A_i^\circ)$、$F=F_1\overset L\otimes_K F_2\in\operatorname{ob}D^b(A^e)$、$R\mathcal{H}\!om_A(E,F)\in\operatorname{ob}D^b(A^\circ)$，其中 $E=E_1\overset L\otimes_K E_2$。则有\Term{commutative-diagram}{交换方块}
\begin{equation}
\resizebox{.94\linewidth}{!}{\begin{tikzcd}[ampersand replacement=\&,column sep=large,row sep=large]
(R\mathcal{H}\!om_{A_1}(E_1,F_1)\overset L\otimes_{A_1}E_1)\overset L\otimes_K(R\mathcal{H}\!om_{A_2}(E_2,F_2)\overset L\otimes_{A_2}E_2)\arrow[r,"(1)"]\arrow[d,"(2)"'] \& (F_1\overset L\otimes_{A_1^e}A_1)\overset L\otimes_K(F_2\overset L\otimes_{A_2^e}A_2)\arrow[d,"(3)"]\\
R\mathcal{H}\!om_A(E,F)\overset L\otimes_A E\arrow[r] \& F\overset L\otimes_{A^e}A,
\end{tikzcd}}
\tag{5.12.1}\end{equation}
其中 (1) 是\Term{arrow}{箭头} 5.3.5 的\Term{tensor-product}{张量积}，(2) 由\Term{tensor-product}{张量积}\Term{arrow}{箭头}
\begin{equation}R\mathcal{H}\!om_{A_1}(E_1,F_1)\overset L\otimes_K R\mathcal{H}\!om_{A_2}(E_2,F_2)\longrightarrow R\mathcal{H}\!om_A(E,F)\tag{5.12.2}\end{equation}
定义，(3) 是显然的\Term{canonical}{典范}\Term{isomorphism}{同构}。验证本质上是显然的。同样可验证，对 $E_i\in\operatorname{ob}D^-(A_i)$、$F_i\in\operatorname{ob}D^+(A_i^e)$、$G_i\in\operatorname{ob}D^-(A_i)$，若 $F_i\overset L\otimes_{A_i}G_i\in\operatorname{ob}D^b(A_i)$、$R\mathcal{H}\!om_{A_i}(E_i,F_i)\in\operatorname{ob}D^b(A_i^\circ)$、$R\mathcal{H}\!om_{A_i}(E_i,F_i\overset L\otimes_{A_i}G_i)\in\operatorname{ob}D^b(K)$、$F\overset L\otimes_A G\in\operatorname{ob}D^b(A)$（其中 $G=G_1\overset L\otimes_K G_2$），则有\Term{commutative-diagram}{交换方块}
\SourcePage{192}
\begin{equation}
\resizebox{.94\linewidth}{!}{\begin{tikzcd}[ampersand replacement=\&,column sep=large,row sep=large]
(R\mathcal{H}\!om_{A_1}(E_1,F_1)\overset L\otimes_{A_1}G_1)\overset L\otimes_K(R\mathcal{H}\!om_{A_2}(E_2,F_2)\overset L\otimes_{A_2}G_2)\arrow[r,"(1)"]\arrow[d,"(2)"'] \& R\mathcal{H}\!om_{A_1}(E_1,F_1\overset L\otimes_{A_1}G_1)\overset L\otimes_K R\mathcal{H}\!om_{A_2}(E_2,F_2\overset L\otimes_{A_2}G_2)\arrow[d,"(3)"]\\
R\mathcal{H}\!om_A(E,F)\overset L\otimes_A G\arrow[r,"5.4.2"'] \& R\mathcal{H}\!om_A(E,F\overset L\otimes_A G),
\end{tikzcd}}
\tag{5.12.3}\end{equation}
其中 (2)、(3) 是 5.12.2 型的\Term{tensor-product}{张量积}\Term{arrow}{箭头}，(1) 是\Term{arrow}{箭头} 5.4.2 的\Term{tensor-product}{张量积}。由方块 5.12.1 和 5.12.3 的\Term{commutative-diagram}{交换性}\Term{pushout}{推出}，对 $E_i\in\operatorname{ob}D(A_i)_{\mathrm{parf}}$、$F_i\in\operatorname{ob}D^b(A_i^e)$，若 $F=F_1\overset L\otimes_K F_2\in\operatorname{ob}D^b(A^e)$，则有\Term{commutative-diagram}{交换方块}
\begin{equation}
\resizebox{.94\linewidth}{!}{\begin{tikzcd}[ampersand replacement=\&,column sep=large,row sep=large]
R\mathcal{H}\!om_{A_1}(E_1,F_1\overset L\otimes_{A_1}E_1)\overset L\otimes_K R\mathcal{H}\!om_{A_2}(E_2,F_2\overset L\otimes_{A_2}E_2)\arrow[r,"(1)"]\arrow[d,"(2)"'] \& (F_1\overset L\otimes_{A_1^e}A_1)\overset L\otimes_K(F_2\overset L\otimes_{A_2^e}A_2)\arrow[d,"(3)"]\\
R\mathcal{H}\!om_A(E,F\overset L\otimes_A E)\arrow[r,"5.6.1"'] \& F\overset L\otimes_{A^e}A,
\end{tikzcd}}
\tag{5.12.4}\end{equation}
其中 (1) 是\Term{arrow}{箭头} $\operatorname{Tr}_{A_i}$（5.6.1）的\Term{tensor-product}{张量积}，(2) 由 5.12.2 定义，(3) 是 5.12.1 中出现的\Term{canonical}{典范}\Term{isomorphism}{同构}。因此，对 $u_i\in\operatorname{Hom}_{A_i}(E_i,F_i\overset L\otimes_{A_i}E_i)$ 有
\begin{equation}\operatorname{Tr}_A(u_1\overset L\otimes_K u_2)=\operatorname{Tr}_{A_1}(u_1)\operatorname{Tr}_{A_2}(u_2),\tag{5.12.5}\end{equation}
其中右边的乘积由\Term{canonical}{典范}\Term{arrow}{箭头}
\[H^0(T,F_1\overset L\otimes_{A_1^e}A_1)\otimes H^0(T,F_2\overset L\otimes_{A_2^e}A_2)\longrightarrow H^0(T,F\overset L\otimes_{A^e}A)\]
定义。

\SGAsection{5.13}{半线性自同态的迹}
设 $A$ 是一个 $K$-\Term{algebra}{代数}，$\varphi$ 是 $A$ 的一个\Term{endomorphism}{自同态}。若 $E,F$ 是左 $A$-\Term{module}{模}，回顾，一个加法\Term{homomorphism}{同态} $u:E\to F$ 称为 $\varphi$-$A$-\Term{semilinear}{线性}，若对任意 $a\in A$、$x\in E$ 都有 $u(ax)=\varphi(a)u(x)$；也就是说，$u$ 是从 $E$ 到由 $F$ 沿 $\varphi$ \Term{restriction-of-scalars}{限制标量}得到的 $A$-\Term{module}{模} ${}_{(\varphi)}F$ 的 $A$-\Term{linear}{线性}\Term{homomorphism}{同态}。若 $F$ 是一个 $A$-\Term{bimodule}{双模}，
\SourcePage{193}
$G$ 是左 $A$-\Term{module}{模}，则有显然的\Term{isomorphism}{同构} ${}_{(\varphi)}F\otimes_A G\xrightarrow\sim{}_{(\varphi)}(F\otimes_A G)$；它（平凡地）\Term{derive}{导出}为 ${}_{(\varphi)}F\overset L\otimes_A G\xrightarrow\sim{}_{(\varphi)}(F\overset L\otimes_A G)$。

假设 $A$ 在 $K$ 上\Term{flat}{平坦}。于是由 5.4.2 得到\Term{arrow}{箭头}
\begin{equation}R\mathcal{H}\!om_A(E,{}_{(\varphi)}F)\overset L\otimes_A G\longrightarrow R\mathcal{H}\!om_A(E,{}_{(\varphi)}(F\overset L\otimes_A G)),\tag{5.13.1}\end{equation}
其中 $E\in\operatorname{ob}D^-(A)$、$F\in\operatorname{ob}D^+(A^e)$、$G\in\operatorname{ob}D^-(A)$，且 $F\overset L\otimes_A G\in\operatorname{ob}D^+(A)$。对 $E\in\operatorname{ob}D(A)_{\mathrm{parf}}$，5.13.1 是\Term{isomorphism}{同构}，因而可由 5.3.5 定义\Term{trace}{迹}\Term{arrow}{箭头}
\begin{equation}\operatorname{Tr}_{A,\varphi}:R\mathcal{H}\!om_A(E,{}_{(\varphi)}(F\overset L\otimes_A E))\longrightarrow{}_{(\varphi)}F\overset L\otimes_{A^e}A.\tag{5.13.2}\end{equation}
特别取 $F=A$，得到\Term{arrow}{箭头}
\begin{equation}\operatorname{Tr}_{A,\varphi}:\operatorname{Hom}_A(E,{}_{(\varphi)}E)\longrightarrow H^0(T,A_{\natural,\varphi}),\tag{5.13.3}\end{equation}
其中 $A_{\natural,\varphi}={}_{(\varphi)}A\otimes_{A^e}A$，是 $A$ 对由形如 $ab-\varphi(b)a$ 的\Term{local-section}{局部截面}\Term{local-scheme}{局部}生成的\Term{submodule}{子模}的商 $K$-\Term{module}{模}。在 $A=K[G]$、$G$ 是\Term{discrete}{离散}群且 $\varphi$ 由 $G$ 的一个\Term{endomorphism}{自同态} $\varphi$ 定义的情形，$A_{\natural,\varphi}$ 是以 $G_{\natural,\varphi}$ 为基的自由 $K$-\Term{module}{模}；其中 $G_{\natural,\varphi}$ 是 $G$ 对\Term{action}{作用} $(g,x)\mapsto\varphi(g)xg^{-1}$ 的商（其\Term{orbit}{轨道}有时称为 $G$ 的“$\varphi$-\Term{conjugacy-class}{共轭类}”）。当 $T$ 是点\Term{topos}{拓扑斯}时，\Term{trace}{迹} 5.14.2 与 Grothendieck 所考虑、并在第 XII 讲中论及的\Term{trace}{迹}一致。

在 5.10 的情形，若 $\varphi,\psi$ 分别是 $A,B$ 的\Term{endomorphism}{自同态}，且与 $A\to B$ 相容，则有唯一的\Term{arrow}{箭头} $\operatorname{Tr}_A:B_{\natural,\psi}\to A_{\natural,\varphi}$ 使以下方块（类似于 5.10.2）\Term{commutative-diagram}{交换}：
\begin{equation}\begin{tikzcd}[column sep=large,row sep=large]
{}_{(\psi)}B=\mathcal{H}\!om_B(B,{}_{(\psi)}B)\arrow[r,"{\operatorname{Tr}_{B,\psi}}"]\arrow[d] & B_{\natural,\psi}\arrow[d,"\operatorname{Tr}_A"]\\
\mathcal{H}\!om_A(B,{}_{(\varphi)}A)\otimes_A B\arrow[r,"\widetilde{\mathrm{ev}}"'] & A_{\natural,\varphi},
\end{tikzcd}\tag{5.13.4}\end{equation}
其中 $\widetilde{\mathrm{ev}}$ 是\Term{evaluation-map}{求值箭头}，左侧竖直\Term{arrow}{箭头}是\Term{restriction-of-scalars}{标量限制}。将 $A\otimes_K M$（相应地
\SourcePage{194}
$B\otimes_K M$）换成 ${}_{(\varphi)}A\otimes_K M$（相应地 ${}_{(\psi)}B\otimes_K M$），可定义类似于 5.10.5 的\Term{arrow}{箭头}，并有类似于 5.10.6、5.10.9、5.10.11 的相容性，其表述留给读者。

设 $G$ 为\Term{finite-group}{有限群}，$\varphi$ 为 $G$ 的一个\Term{endomorphism}{自同态}。对 $g\in G$，以 $Z_{g,\varphi}$ 表示 $g$ 在 $G$ 中的 $\varphi$-\Term{centralizer}{中心化子}，即满足 $\varphi(x)gx^{-1}=g$ 的 $x\in G$ 全体。置 $A=K[G]$，设 $E\in\operatorname{ob}D(A)_{\mathrm{parf}}$、$N\in\operatorname{ob}D^b(K)$。由上述相容性容易\Term{pushout}{推出}，对 $u\in\operatorname{Hom}_A(E,{}_{(\varphi)}(M\overset L\otimes_K E))$、$g\in G$，有
\begin{equation}\operatorname{Tr}_K(gu)=|Z_{g^{-1},\varphi}|\operatorname{Tr}_{A,\varphi}(u)(g^{-1})\tag{5.13.5}\end{equation}
（验证的细节留给读者）。
